% !TeX root = ../教材题库_全部条目_转写工作稿.tex
% Source: XJ-B1-C06-HW-P257-Y0449
% 第6章 / 6.4 用样本估计总体 / 教材 p.257 / PDF 第264页
% OCR未校对：公式与文字均可能有误；未完成数学TeX转写；禁止直接发给学生
% 原题及相邻原解截图（核对用）：../题目原图/XJ-B1-C06-HW-P257-Y0449-01.webp
% 本文件只在转写完成后作为题干使用；OCR 条目尚不是可靠题干。
\begin{bankitem}{习题6.4 第13题}{第6章 / 6.4 用样本估计总体 / 教材 p.257 / PDF 第264页}{XJ-B1-C06-HW-P257-Y0449}{OCR未校对：公式与文字均可能有误；未完成数学TeX转写；禁止直接发给学生}
\begin{ocrdraft}
[教材p.257 / PDF第264页]
\par
13．从一家企业生产的某种产品中抽取100件，测量这些产品的一项质量指标
\par
值，由测量结果得到如下频数分布表：
\par
质量指标值分组[75，85）
\par
[85，95）
\par
[95，105）
\par
[105，115)[115，125）
\par
6
\par
26
\par
38
\par
22
\par
8
\par
（1)根据以上数据作出频率分布直方图；
\par
（2）估计这种产品质量指标值的平均数及方差（同一组中的数据用该组区间的
\par
组中值作代表）；
\par
（3）根据以上抽样调查数据，能否认为该企业生产的这种产品符合“质量指标
\par
值不低于95的产品至少要占全部产品的80％”的规定？
\end{ocrdraft}
\end{bankitem}
